% answers_short.tex —— 2016 年·北京卷·理科 · 简短答案速查
% 由 tools/make_short_answers.py 自 answers.tex 的 \gksolhead/\gkanswer 提取，
% 并与 2016/2016_数学真题_北京理_答案要点.md 逐题比对后生成，勿手改（改源 MD 后重跑）。

\begin{gknotice}
\textbf{使用说明}：本节只给最终答案便于快速核对，完整推导见解析版。答案为编写者独立推导，\textbf{非官方标准答案}。
\end{gknotice}

\gkpart{选择题}
\begin{center}\begin{tabular}{ccccccccc}
\toprule
题号 & 1 & 2 & 3 & 4 & 5 & 6 & 7 & 8 \\
\midrule
答案 & C & C & B & D & C & A & A & B \\
\bottomrule
\end{tabular}\end{center}

\gkpart{填空题}
\begin{itemize}[leftmargin=2.6em]
  \item[\textbf{9.}] \(-1\)
  \item[\textbf{10.}] \(60\)
  \item[\textbf{11.}] \(2\)
  \item[\textbf{12.}] \(6\)
  \item[\textbf{13.}] \(2\)
  \item[\textbf{14.}] ① \(2\)；② \((-\infty,-1)\)。
\end{itemize}

\gkpart{解答题}
\begin{itemize}[leftmargin=2.6em]
  \item[\textbf{15.}] (1) \(B=\dfrac{\pi}{4}\)；(2) 最大值为 \(1\)
  \item[\textbf{16.}] (1) \(40\) 人；(2) \(\dfrac38\)；(3) \(\mu_0>\mu_1\)
  \item[\textbf{17.}] (1) 见解析版；(2) \(\dfrac{\sqrt3}{3}\)；(3) 存在，\(\dfrac{AM}{AP}=\dfrac14\)
  \item[\textbf{18.}] (1) \(a=2\)，\(b=\mathrm e\)；(2) 单调递增区间为 \((-\infty,+\infty)\)，无递减区间
  \item[\textbf{19.}] (1) \(\dfrac{x^2}{4}+y^2=1\)；(2) 见解析版
  \item[\textbf{20.}] (1) \(G(A)=\{2,5\}\)；(2) 见解析版；(3) 见解析版
\end{itemize}
